\documentclass[11pt]{article}
\usepackage[margin=0.85in]{geometry}
\usepackage[T1]{fontenc}
\usepackage{amsmath,amssymb,booktabs,graphicx,xurl,subcaption,float}
\usepackage[hidelinks]{hyperref}
\newcommand{\widerow}[1]{\noindent\makebox[\linewidth][c]{#1}}
\graphicspath{{report_figs/}{outputs/task1/}{outputs/task2/}{outputs/task3/}{outputs/bonus/}}

\title{CS 498 HW0 Report\\Perspective, Homography, and Epipolar Geometry}
\author{Siddharth Dinkar (sdinkar2)}
\date{Fall 2026}

\begin{document}
\maketitle

\section{Homography and 2D insertion}

\paragraph{Normalization.} Before building any linear system I precondition each 2D point
set separately. For a set with centroid $\mathbf c$ and mean distance $d$ from
that centroid, the similarity
\[
T=\begin{bmatrix} s & 0 & -s\,c_x\\ 0 & s & -s\,c_y\\ 0 & 0 & 1\end{bmatrix},
\qquad s=\sqrt{2}/d ,
\]
moves the centroid to the origin and scales the mean radius to $\sqrt2$. I build
$T_{xy}$ from the court points and $T_{uv}$ from the pixels. Without this step
the columns of the design matrix range from $1$ to $u\cdot x\approx 10^4$ and the
smallest singular vector is poorly determined.

\paragraph{DLT.} With $\mathbf p=[x,y,1]^\top$ and $H$ written as rows
$\mathbf h_1,\mathbf h_2,\mathbf h_3$, the relation $\lambda[u,v,1]^\top=H\mathbf p$
loses $\lambda$ by substituting the third component into the first two, giving
two rows per point,
\[
\begin{bmatrix}-\mathbf p^\top & \mathbf 0^\top & u\,\mathbf p^\top\\
\mathbf 0^\top & -\mathbf p^\top & v\,\mathbf p^\top\end{bmatrix}
\begin{bmatrix}\mathbf h_1\\ \mathbf h_2\\ \mathbf h_3\end{bmatrix}=\mathbf 0 .
\]
Stacking all points gives a $2N\times 9$ matrix $A$. The right singular vector of
$A$ with the smallest singular value is the unit vector minimizing
$\|A\mathbf h\|$; I reshape it row-major to $3\times3$, undo the normalization
with $H=T_{uv}^{-1}H'T_{xy}$, and divide by the Frobenius norm.

\paragraph{Logo placement.} A logo pixel $(p_x,p_y)$ in an $h\times w$ image is first mapped
to the metric rectangle with lower-left corner $(x_0,y_0)$ and size $(W,L)$,
then through $H$:
\[
M_{\text{logo}\to\text{court}}=\begin{bmatrix} W/w & 0 & x_0\\ 0 & -L/h & y_0+L\\ 0&0&1\end{bmatrix},
\qquad
M_{\text{logo}\to\text{image}}=H\,M_{\text{logo}\to\text{court}} .
\]
The negative entry flips the vertical axis because pixel rows grow downward
while court $y$ grows upward. The provided bilinear inverse warp uses the
composed matrix. Compositing is the handout equation applied per pixel and
channel, $I_\text{out}=\alpha I_\text{logo}+(1-\alpha)I_\text{court}$, where
$\alpha$ is the fourth channel of the warped logo.

\paragraph{Results.} The maximum corner-transfer error on the four basketball corners is
$5.1\times10^{-13}$~px, which is zero up to floating point. Four points give
eight equations for eight unknowns, so the fit is exact by construction.

\begin{figure}[h]\centering
\widerow{%
\begin{subfigure}{0.585\linewidth}\includegraphics[width=\linewidth]{task1a_basketball_homography.png}
\caption{1A: given corners as crosses, $H$-projected corners as circles.}\end{subfigure}\hspace{0.01\linewidth}%
\begin{subfigure}{0.585\linewidth}\includegraphics[width=\linewidth]{task1b_basketball_insertion.png}
\caption{1B: $6\times3$~m logo composited at lower-left $(23.0,\,2.5)$~m.}\end{subfigure}}
\caption{Basketball homography fit and logo insertion.}\label{fig:t1}
\end{figure}

\subsection{Tennis transfer and camera height}

The course frame puts the origin at the center of the near baseline with $+x$
across the court and $+y$ toward the far baseline. Using the regulation length
$23.77$~m, doubles width $10.97$~m and singles width $8.23$~m, the metric
coordinates for the supplied anchors are:

\begin{center}\small
\begin{tabular}{@{}llrr@{}}\toprule
 & intersection & $x$ (m) & $y$ (m)\\\midrule
A & near left doubles corner & $-5.485$ & $0$\\
B & near right doubles corner & $5.485$ & $0$\\
C & far left doubles corner & $-5.485$ & $23.77$\\
D & far right doubles corner & $5.485$ & $23.77$\\
E & near left singles corner & $-4.115$ & $0$\\
F & near right singles corner & $4.115$ & $0$\\
G & far left singles corner & $-4.115$ & $23.77$\\
H & far right singles corner & $4.115$ & $23.77$\\\bottomrule
\end{tabular}
\end{center}

Running the same three functions on frame 0 of \path{match005.mp4} gives an
eight-anchor reprojection RMSE of $3.67$~px (Fig.~\ref{fig:t1c}). The near
anchors carry 4 to 6~px and the far anchors about 1~px.

\begin{figure}[h]\centering
\begin{subfigure}{0.49\linewidth}\includegraphics[width=\linewidth]{task1c_tennis_correspondences.png}
\caption{Supplied anchors A--H.}\end{subfigure}\hfill
\begin{subfigure}{0.49\linewidth}\includegraphics[width=\linewidth]{task1c_tennis_alignment.png}
\caption{Given anchors in green, reprojected in orange.}\end{subfigure}\\[4pt]
\begin{subfigure}{0.49\linewidth}\includegraphics[width=\linewidth]{task1c_tennis_insertion.png}
\caption{Logo in $[-3,3]\times[0.5,3.5]$~m.}\end{subfigure}
\caption{Tennis transfer using the coordinates above and unchanged Task 1 code.}\label{fig:t1c}
\end{figure}

\paragraph{Camera height.} For a planar scene the camera matrix loses its third column, so
$H\sim K[\mathbf r_1\;\mathbf r_2\;\mathbf t]$; $\mathbf r_3$ never appears
because every court point has $Z=0$. With $K$ unknown this product cannot be
split apart. A floor point projects to $u=f\,X_c/Z_c+c_x$, so moving the camera
twice as far away doubles $Z_c$ and doubling the focal length doubles $f$; the
two changes cancel and every floor pixel stays put. Two cameras at different
heights therefore produce the same $H$. Counting degrees of freedom says the
same thing: $H$ has 8, while $K$ (5), $R$ (3) and $\mathbf t$ (3) have 11, so a
three-parameter family of cameras shares each $H$. The basketball scene breaks
the tie with the four backboard corners at $Z>0$. Those points constrain the
column of $P$ that the plane discards, so the full $3\times4$ matrix, and from
it $K$, $R$ and the height, become observable without knowing $K$ in advance.

\section{Camera projection and 3D insertion}

\paragraph{Projection DLT.} The construction is the Task 1 DLT with a 3D point. With
$\mathbf X=[X,Y,Z,1]^\top$ and $P$ written as rows $\mathbf p_1,\mathbf p_2,\mathbf p_3\in\mathbb R^4$,
each landmark contributes the rows
$[-\mathbf X^\top,\;\mathbf 0^\top,\;u\mathbf X^\top]$ and
$[\mathbf 0^\top,\;-\mathbf X^\top,\;v\mathbf X^\top]$ to a $2N\times12$ system.
I normalize the 3D points with a $4\times4$ similarity that scales the mean
radius to $\sqrt3$ and the pixels with the $3\times3$ transform from Task 1,
take the last row of $V^\top$, reshape to $3\times4$, undo with
$P=T_{uv}^{-1}P'T_{xyz}$, and fix the Frobenius norm. $P$ has 11 degrees of
freedom so at least 6 points are needed; all eight landmarks are used.

\paragraph{Role of $K$.} If every landmark had $Z=0$, the four $Z$ columns of the $2N\times12$
design matrix would be identically zero. The three entries of $P$ in its $Z$
column would then multiply nothing, so any values for them satisfy
$A\mathbf p=\mathbf 0$ and the null space has dimension at least four; the SVD
returns an arbitrary vector from it and the problem collapses to the homography
of Task 1. Whether that is a problem depends on $K$. If $K$ is known, the plane
is enough: $K^{-1}H=[\mathbf r_1\;\mathbf r_2\;\mathbf t]$ up to scale, and
$\mathbf r_3=\mathbf r_1\times\mathbf r_2$ completes the rotation. If $K$ is
unknown, as here, $H$ cannot be factored and points off the plane are needed to
pin down all four columns of $P$. Deleting the $Z$ column of my estimated $P$
gives back my Task 1 $H$ to within $0.0014$ after scaling both to unit norm,
which confirms the two tasks estimate the same camera.

\paragraph{Results.} The mean and maximum landmark reprojection errors are $2.75$~px and
$5.10$~px (Fig.~\ref{fig:t2}). The court corners are within $1.7$~px; the four
small, distant backboard corners carry 4 to 5~px. The provided decomposition
returns
\[
K\approx\begin{bmatrix}-763.9&-258.9&481.8\\0&-361.8&-202.2\\0&0&1\end{bmatrix},\quad
R\approx\begin{bmatrix}-0.706&-0.708&-0.019\\0.602&-0.614&0.510\\-0.372&0.349&0.860\end{bmatrix},\quad
\mathbf t\approx\begin{bmatrix}22.72\\-39.86\\17.21\end{bmatrix},
\]
with $C=(46.45,\,-14.41,\,5.95)$~m, so the optical center sits $5.95$~m above the
court. Loading \path{task2_results.npz} and computing $-R^\top\mathbf t$
reproduces $C$ to $7\times10^{-15}$~m per component; $R^\top R=I$ with
$\det R=1$; and $P\,[C;1]=\mathbf 0$, so $C$ is the null vector of $P$. Following
the handout, $K$, $R$ and $\mathbf t$ are treated as noisy DLT estimates rather
than a calibration; the negative focal terms and large skew in $K$ reflect
that.

\begin{figure}[h]\centering
\begin{subfigure}{0.49\linewidth}\includegraphics[width=\linewidth]{task2a_reprojection.png}
\caption{2A: eight landmarks reprojected by $P$.}\end{subfigure}\hfill
\begin{subfigure}{0.49\linewidth}\includegraphics[width=\linewidth]{task2b_bunny_insertion.png}
\caption{2B: bunny at center court, colored by height.}\end{subfigure}\\[4pt]
\begin{subfigure}{0.49\linewidth}\includegraphics[width=\linewidth]{task2c_camera_pose.png}
\caption{2C: recovered camera center and viewing direction.}\end{subfigure}
\caption{Camera matrix estimation, 3D insertion, and pose.}\label{fig:t2}
\end{figure}

\section{Two-view epipolar geometry}

\paragraph{Method.} Each match row is $(u_1,v_1,u_2,v_2)$. The two views are normalized
independently with $T_1$ and $T_2$ as in Task 1. The epipolar constraint
$\mathbf x_2^\top F\mathbf x_1=0$ is already linear in $F$, so each match gives a
single row
$[u_2u_1,\;u_2v_1,\;u_2,\;v_2u_1,\;v_2v_1,\;v_2,\;u_1,\;v_1,\;1]$ of an
$N\times9$ matrix. The last row of $V^\top$ reshaped to $3\times3$ is $\tilde F$.
A second SVD enforces rank two, $F'=U\,\mathrm{diag}(\sigma_1,\sigma_2,0)\,V^\top$,
which is the nearest rank-2 matrix in Frobenius norm. Because the constraint is
bilinear the denormalization is $F=T_2^\top F' T_1$, a transpose rather than an
inverse. $F$ is then scaled to unit Frobenius norm.

\paragraph{Results.} From the eight supplied matches,
\[
F\approx\begin{bmatrix}-2.20\times10^{-9}&1.06\times10^{-7}&-1.38\times10^{-4}\\
2.68\times10^{-7}&2.54\times10^{-8}&7.05\times10^{-4}\\
-2.97\times10^{-4}&-1.46\times10^{-3}&1.00\end{bmatrix},
\qquad \operatorname{rank}F=2 .
\]
Over all 200 evaluation matches the symmetric epipolar distance has mean
$5.58$~px and median $4.52$~px (Fig.~\ref{fig:t3}).

\paragraph{Why the residuals are nonzero.} Four effects keep the residuals away from zero. The 200 keypoints come from a
feature detector on real photographs and are localized to roughly a pixel, so
each match carries noise before $F$ is estimated. $F$ was fitted to only 8 of
the 200 matches, and the other 192 cannot be predicted better than the fit
itself. The rank-2 projection moves $F$ off the exact null vector, so even the 8
fitted points no longer lie exactly on their lines; their mean error is
$1.7$~px. Finally, $F$ assumes ideal pinhole cameras and cannot absorb lens
distortion. Refitting with all 200 matches lowers the mean to $2.8$~px, which
shows the remaining error is a property of the data rather than the code.

\begin{figure}[h]\centering
\widerow{%
\begin{subfigure}{0.585\linewidth}\includegraphics[width=\linewidth]{task3b_correspondences.png}
\caption{Ten of the supplied correspondences.}\end{subfigure}\hspace{0.01\linewidth}%
\begin{subfigure}{0.585\linewidth}\includegraphics[width=\linewidth]{task3b_epipolar_lines.png}
\caption{Epipolar lines from $F$; each passes near its match.}\end{subfigure}}
\caption{Two-view geometry on the fountain pair.}\label{fig:t3}
\end{figure}

\section{Limitations and improvements}

Annotation precision. The basketball error is exactly zero only because four
points give eight equations for eight unknowns; any click error is absorbed into
$H$ instead of appearing as residual, so that number says nothing about
accuracy. The tennis ($3.7$~px), landmark ($2.8$~px) and epipolar ($5.6$~px)
fits are overdetermined, and their residuals are the click and detection noise
finally becoming visible. The near-baseline tennis anchors are the worst
because perspective spreads a pixel of error there over more meters of court.

Degenerate configurations. The homography fails if three or more of the
points are collinear. The projection DLT fails when all landmarks are coplanar,
the $Z=0$ case discussed in Section 2. The eight-point algorithm fails for a
pure rotation between views, since without a baseline there are no epipolar
lines, and when all matched points lie on one plane, since a homography then
explains the matches and $F$ is not unique.

Improvements. First, the DLT minimizes an algebraic residual rather than pixel
distance; a few iterations of nonlinear least squares on reprojection or
Sampson error after the linear solve would reduce the geometric error directly.
Second, $F$ should be estimated from all 200 matches with RANSAC rather than
from eight hand-picked ones, so that a large inlier set averages out the
detection noise. Third, a radial distortion model for the broadcast tennis
frames would remove the systematic near-versus-far error pattern.

\section*{Optional foreground-preserving video bonus}

I implemented the complete bonus pipeline in \path{bonus_tennis_video.py}
without any provided helper. Video loading, background estimation, the
foreground mask, logo insertion, composition and output writing are all in that
file; the only reused code is my Task 1 \texttt{estimate\_homography},
\texttt{logo\_to\_image\_homography} and \texttt{alpha\_blend}.

\paragraph{Background estimate.} Every clip has a fixed camera, so the per-pixel temporal
median of 45 frames sampled evenly across the clip removes the moving players
and leaves the empty court.

\paragraph{Foreground alpha.} For each frame I take the maximum absolute color difference
to that background, pass it through a smooth ramp between $0.10$ and $0.28$,
clean it with a morphological opening (speckle) and closing (holes inside
players), dilate slightly, and Gaussian blur to feather the edges. The result is
a soft mask $m\in[0,1]$.

\paragraph{Composition.} The logo is warped with the Task 1 homography built from the clip's
eight supplied image anchors and the metric A--H coordinates above, dimmed by
the ratio of live to background luminance so shadows fall across it, and blended
with effective alpha $\alpha_\text{logo}(1-m)$. Where a player is present the
effective alpha drops to zero and the original frame pixel remains on top; the
player is never cut out and pasted back. Compositing runs only inside a box
around the logo footprint. The representative frame for each clip is the one
where $m$ overlaps the logo most (Fig.~\ref{fig:bonus}), and
\path{match005_occlusion_aware.mp4} holds all 250 composited frames.

\begin{figure}[p]\centering
\begin{subfigure}{0.49\linewidth}\includegraphics[width=\linewidth]{match001_representative.png}
\caption{match001, frame 500.}\end{subfigure}\hfill
\begin{subfigure}{0.49\linewidth}\includegraphics[width=\linewidth]{match005_representative.png}
\caption{match005, frame 3.}\end{subfigure}\\[4pt]
\begin{subfigure}{0.49\linewidth}\includegraphics[width=\linewidth]{match028_representative.png}
\caption{match028, frame 131.}\end{subfigure}\hfill
\begin{subfigure}{0.49\linewidth}\includegraphics[width=\linewidth]{match091_representative.png}
\caption{match091, frame 356.}\end{subfigure}\\[4pt]
\begin{subfigure}{0.49\linewidth}\includegraphics[width=\linewidth]{match125_representative.png}
\caption{match125, frame 578.}\end{subfigure}\hfill
\begin{subfigure}{0.49\linewidth}
\includegraphics[width=\linewidth]{match005_mask_crop.png}
\caption{Soft mask for the match005 frame, logo region.}\end{subfigure}
\caption{Occlusion-aware insertion: the player stays in front of the logo in every clip.}\label{fig:bonus}
\end{figure}

\section*{AI-use declaration}
Check one statement.

\begin{itemize}
  \item[$\square$] I did not use generative AI for this homework.
  \item[$\boxtimes$] I did not use generative AI for any required part. I used it only in the optional bonus marked ``Gen-AI leverage allowed.'' Tool/model, purpose, and affected files: Claude (Anthropic), used to design and write the bonus compositor; affected files are \path{bonus_tennis_video.py} and the generated \path{outputs/bonus/} only.
\end{itemize}

I confirm that generative AI did not write or complete my required code,
select required correspondences, derive my required answers, generate my
analysis, or write my required report.

\end{document}
